\section{The geometric mechanism}
\label{sec:geometric-closure}

This section isolates what the analysis has to produce. Throughout,
$0<r_-<r_+$, $\T=\R/(2\pi\mathbb Z)$, and
\begin{equation}\label{eq:intro-metric}
 g=E(r)\,dr^2+r^2\,d\theta^2,\qquad E=1+w^2.
\end{equation}
A function is flat at an endpoint when it extends smoothly there
with every derivative zero. A height has a constant endpoint extension
when its value has a limit and every positive-order derivative of its
difference from that limit vanishes. Extending these functions and
slightly enlarging the annulus then gives actual open collars on
which $w=0$ and the height is constant. Thus endpoint flatness, rather
than a merely small error at the boundary, is the appropriate analytic
target.
\label{def:flat-endpoint}

\subsection{Removing one coordinate from the metric}

An embedding of the form $F=(P,z)$ is isometric exactly when
\[
 P^*g_{\R^2}=\bar g:=g-dz\otimes dz.
\]
We need $\bar g$ to be a positive flat metric and its planar
development to be injective. The local condition is a scalar equation.
For a symmetric tensor $B$, write
$\det_g B=\det(B_{ij})/\det(g_{ij})$.

\begin{lemma}\label{lem:darboux-invariant}
For a smooth height $z$ with $v=|dz|_g^2<1$,
\begin{equation}\label{eq:darboux-curvature}
 K_{g-dz^2}=
 \frac{K_g(1-v)-\det_g\nabla_g^2z}{(1-v)^2}.
\end{equation}
In particular, the positive metric $g-dz^2$ is flat precisely when
\begin{equation}\label{eq:intro-darboux}
 \det_g\nabla_g^2z=K_g(1-|dz|_g^2).
\end{equation}
\end{lemma}

\begin{proof}
In a $g$-orthonormal frame with its first covector parallel to $dz$,
$\bar g$ has matrix $\diag(1-v,1)$. Regard $(x,z(x))$ as a graph
in $(\Sigma,\bar g)\times\R$, with induced metric $g$. Its upward
unit normal has vertical component $\sqrt{1-v}$, and its second
fundamental form satisfies
$\nabla_g^2z=\sqrt{1-v}\,\mathrm{II}$.
The ambient sectional curvature of its tangent plane is
$(1-v)K_{\bar g}$, because horizontal projection multiplies squared
area by $1-v$. The Gauss equation therefore reads
\[
 K_g=(1-v)K_{\bar g}+(1-v)^{-1}\det_g\nabla_g^2z,
\]
which gives the formula.
\end{proof}

This proof also explains the equation: it is exactly the compatibility
condition for the two horizontal coordinates after a vertical
coordinate has been chosen. No local embedding theorem is being
used to bypass the global reconstruction.

\subsection{Why the development is a graph over a round disk}

Flatness on an annulus alone leaves a holonomy problem. Filling the
inner collar by a planar disk removes it. Half-turn symmetry then
ensures that this disk is centered correctly.

\begin{proposition}[Planar reconstruction]\label{lem:disk-development}
Suppose $w$ vanishes on open collars of both ends and $z$ is constant
there, with $z(r,\theta+\pi)=z(r,\theta)$. If $g-dz^2$ is positive
and flat, extend $w=0$ and $z=c_L$ across the inner disk. There is an
orientation-preserving diffeomorphism
$P:D_{r_+}\longrightarrow D_{r_+}$ with
\[
 P^*g_{\R^2}=g-dz^2,
 \qquad P=\id\text{ on the outer collar},
 \qquad P=R_\gamma\text{ on the inner disk}
\]
for some angle $\gamma$. Consequently $(P,z)$ is an isometric
embedded disk.
\end{proposition}

\begin{proof}
The extended metric is smooth on the disk, including its center,
where it is Euclidean. Take a global oriented orthonormal coframe
$\eta_1,\eta_2$, equal to $dx,dy$ on the outer collar, and let $\alpha$
be its connection form, with
$d\eta_1=\alpha\wedge\eta_2$ and
$d\eta_2=-\alpha\wedge\eta_1$. Flatness and simple connectivity
give $\alpha=d\psi$, normalized by $\psi=0$ on that collar. Thus
$e^{i\psi}(\eta_1+i\eta_2)$ is closed. Its primitive $P$, normalized
to be the identity on the collar, is an oriented local isometry to
the plane.

The half-turn $I:x\mapsto-x$ preserves the metric.
Both $P$ and $-P\circ I$ are normalized developments, so they agree:
two oriented planar developments of a connected flat surface differ
by a fixed Euclidean motion, as follows by differentiating their
relative rotation. Hence $P$ is odd. On the inner Euclidean disk it
is a Euclidean motion, and oddness removes its translation. This
gives the asserted rotation there.

Extend $P$ by the identity outside the disk. It is a proper local
diffeomorphism of $\R^2$, hence a covering. It has one sheet outside
a compact set, so it is a diffeomorphism. Since the exterior is
fixed, the disk maps onto itself. The height graph $(P,z)$ is
therefore embedded, with metric $\bar g+dz^2=g$.
\end{proof}

\begin{figure}[htbp]
\centering
\begin{tikzpicture}[scale=.87,>=Stealth,font=\small]
\begin{scope}
\filldraw[fill=blue!6,draw=black] (0,0) circle (1.65);
\filldraw[fill=green!12,draw=green!45!black] (0,0) circle (.58);
\draw[dashed] (0,0) circle (1.42);
\draw[->] (0,0)--(1.25,.55) node[right] {$r$};
\node[below,align=center] at (0,-1.9) {intrinsic disk\\metric $g-dz^2$};
\end{scope}
\draw[->,thick] (2.1,0)--(3.9,0) node[midway,above] {$P$};
\node[align=center] at (3,-.65) {flatness\\and no holonomy};
\begin{scope}[xshift=6cm]
\filldraw[fill=blue!6,draw=black] (0,0) circle (1.65);
\filldraw[fill=green!12,draw=green!45!black] (0,0) circle (.58);
\draw[dashed] (0,0) circle (1.42);
\draw[->] (.1,.1) arc[start angle=45,end angle=210,radius=.25];
\node at (0,.95) {$P=\id$ near the rim};
\node[below,align=center] at (0,-1.9) {planar image\\same round disk};
\end{scope}
\end{tikzpicture}
\caption{The inner fill removes the annular holonomy. Normalization
fixes the outer collar; half-turn symmetry fixes the center of the
inner disk. Properness then upgrades the local development to a
global diffeomorphism. The sketch records these constraints, not
the pointwise distortion of $P$.}
\label{fig:development}
\end{figure}

\subsection{A stationary cap which distinguishes shapes}

We need more than a cap that closes the graph. A completely symmetric
cap could hide the relative rotation. The following explicit marker
reduces the possible congruences to a finite set.

\begin{lemma}\label{lem:fixed-cap}
A smooth graph over $D_{r_+}$ which is horizontal near its boundary
extends to a smooth embedded sphere. The complementary cap may be
chosen to contain the unique global maximum of Gaussian curvature,
with distinct principal curvatures at that point.
\end{lemma}

\begin{proof}
Translate the boundary height to zero and extend the horizontal
collar to a radius $R>r_+$. Choose a smooth nondecreasing function
$\chi:[0,1]\to[0,1]$, constant at its endpoint values near the
endpoints and satisfying $\chi(1-t)=1-\chi(t)$. The unit-speed
meridian starting at $(R,0)$ with
\[
 r'(s)=\cos(\pi\chi(s/L)),\qquad
 z'(s)=-\sin(\pi\chi(s/L)),\qquad 0\le s\le L,
\]
stays at radius at least $R$ and returns to radius $R$ with inward
horizontal tangent. Its height decreases by a positive constant
times $L$. Take $L$ large enough to pass below the minimum of the
graph, and finish with a horizontal disk to the axis. Rotation of
this meridian gives a smooth embedded cap disjoint from the graph
except at its attaching collar.

Place a bump in the unused horizontal annulus $r_+<r<R$. In local
Cartesian coordinates choose $A,B>0$ with $A\ne B$, put
$\xi=Ax^2+By^2$, and set
\[
 u(x,y)=\frac{k}{2}\int_\xi^{\rho^2}
       e^{-s/(\rho^2-s)}\,ds\quad(\xi<\rho^2),
 \qquad u=0\quad(\xi\ge\rho^2).
\]
This is smooth across its support boundary. If
$\Xi=\xi/(\rho^2-\xi)$, the curvature of its graph is
\begin{equation}\label{eq:marker-curvature}
 K=\frac{ABk^2e^{-2\Xi}(1-2\xi\Xi')}
 {\bigl[1+k^2e^{-2\Xi}(A^2x^2+B^2y^2)\bigr]^2}.
\end{equation}
At the center it equals $ABk^2$. Elsewhere its numerator is either
nonpositive or strictly smaller, and its denominator is at least
one. Choose $k$ to make this value larger than the curvature of the
rest of the sphere. Then choose $\rho$ small enough to fit the
support and keep its height, at most $k\rho^2/2$, away from the
bottom cap. The principal curvatures at the center are $-kA,-kB$
for the upward normal, so they are distinct.
\end{proof}

\begin{figure}[htbp]
\centering
\begin{tikzpicture}[x=.8cm,y=.8cm,>=Stealth,font=\small]
\draw[thick,blue!70!black]
 (-3.3,0)--(-2.8,0) .. controls (-2.5,.1) and (-2.5,.9) .. (-1.5,.75)
 .. controls (-1,.65) and (-.8,.65) .. (-.5,.65)--(.5,.65)
 .. controls (.8,.65) and (1,.65) .. (1.5,.75)
 .. controls (2.5,.9) and (2.5,.1) .. (2.8,0)--(3.3,0);
\draw[thick,gray!75] (3.3,0)--(4,0)
 .. controls (5.5,0) and (5.5,-2) .. (4,-2)--(-4,-2)
 .. controls (-5.5,-2) and (-5.5,0) .. (-4,0)--(-3.3,0);
\fill[red!75!black] (3.7,0) circle (2.5pt);
\draw[dashed,gray] (-2.8,0)--(-2.8,-2);
\draw[dashed,gray] (2.8,0)--(2.8,-2);
\node[above] at (0,1.4) {constant inner disk};
\draw[->] (0,1.3)--(0,.7);
\node[above] at (-3.3,.3) {flat collar};
\node[above] at (3.8,.65) {marker};
\draw[->,red!75!black] (3.8,.6)--(3.7,.1);
\node[below] at (0,-2.1) {fixed cap below the entire moving graph};
\end{tikzpicture}
\caption{A meridional sketch of the closure. The actual blue graph
depends on angle; its whole rotation stays inside the central
cylinder and above the bottom disk. The marker lies in the
stationary horizontal extension.}
\label{fig:geometric-closure}
\end{figure}

\subsection{The annulus implies the theorem}

\begin{theorem}[Geometric closure]\label{thm:geometric-closure}
If a nonradial, half-turn invariant height $z$ for
\eqref{eq:intro-metric} is constant on two collars where $w=0$, and
$g-dz^2$ is positive and flat, then all the conclusions of
Theorem~\ref{thm:main} hold.
\end{theorem}

\begin{proof}
Develop the filled disk by Proposition~\ref{lem:disk-development}
and attach the cap of Lemma~\ref{lem:fixed-cap}. On the disk set
$F_\tau=R_{-2\pi\tau}\circ F\circ R_{2\pi\tau}$, and leave the
cap fixed. The intrinsic rotation preserves $g$. On the inner disk
and outer collar its action cancels the ambient rotation, since
$P$ is a centered rotation there and $z$ is constant. The formula
therefore glues smoothly in both variables and is one-periodic.
Each disk image is a rotation of a graph over the same round disk;
the separated cap closes it without introducing intersections.

The disk images cannot all coincide. Otherwise
$F^{-1}R_\vartheta F$ would be an intrinsic isometry equal to
$R_\vartheta$ on the outer collar. An isometry of a connected
Riemannian manifold is determined by its value and derivative at an
interior point, by geodesic uniqueness. Thus it would equal
$R_\vartheta$ everywhere. Its height component would force
$z(r,\theta+\vartheta)=z(r,\theta)$ for every $\vartheta$, contrary
to nonradiality.

Any congruence between two closed images fixes the unique curvature
maximum in the cap. It preserves the outward normal and the two
distinct principal lines there. There are at most four orthogonal
maps with these properties. If all the closed images were congruent,
their continuous family would therefore lie in a finite set and
would be constant. Its disk portion would be constant too, which
we have just excluded.

Finally, the moving disk is an ambient rotation of itself, while
the rest of the surface is stationary. This proves invariance of
area and of the oriented mean-curvature integral. If $Z_\tau$ is
the height over the planar disk and $c_B$ is the bottom height, the
part of the enclosed volume inside its cylinder is
\[
 \int_{D_{r_+}}(Z_\tau(y)-c_B)\,dy
 =\int_{D_{r_+}}(Z_0(R_{2\pi\tau}y)-c_B)\,dy,
\]
which is constant by rotation invariance of planar area. The
remaining volume is stationary.
\end{proof}

The task left for the analysis is now precise: construct one
nonradial annular height with a positive flat remainder metric and
constant endpoint extensions. The family itself then follows
without an additional existence argument.
